\subsection{维数与截切序列}

设 A 为诺特环，M 为有限生成 A-模，S 为 A 的根基的一个子集，S 为 A 中由 S 生成的理想，而 a 为 M 的湮灭子。我们有

\[
\dim_ {\mathrm{A}} (\mathrm{M}) = \dim (\mathrm{A} / \mathrm{a}) \leqslant \dim (\mathrm{A});\tag{6}
\]

因此，若 A 是局部环，则有 \(\dim_{\mathbf{A}}(\mathbf{M}) < +\infty\)。此外，根据 II, § 4, no. 4 中 Prop. 18 的推论，A-模 M/SM 的支集等于 V(a + S)，因而

\[
\dim_ {A} (M / S M) = \dim (A / (a + \mathfrak {S}))  .\tag{7}
\]

当 S 有限时，或者 M 不为 0 时，我们有不等式

\[
\dim_ {\mathbf {A}} (\mathbf {M} / \mathrm{SM}) \leqslant \dim_ {\mathbf {A}} (\mathbf {M}) \leqslant \operatorname{Card} (\mathbf {S}) + \dim_ {\mathbf {A}} (\mathbf {M} / \mathrm{SM});\tag{8}
\]

当 S 有限时，这由 no. 1 的 prop. 2 以及上面的公式 (6) 和 (7) 得出，而 S 无限的情形是平凡的。

定义 1. --- 设 A 为诺特局部环，M 为非零有限生成 A-模，S 为 A 的极大理想 \(\mathfrak{m}_{\mathrm{A}}\) 的子集。若有

\[
\dim_ {A} (M) = \operatorname{Card} (S) + \dim_ {A} (M / S M).\tag{9}
\]

若 S 是 M 的截切集，则有 \(\operatorname{Card}(S) \leqslant \dim_{\mathbf{A}}(\mathbf{M})\)，所以 S 是有限集。若族 \((x_i)_{i \in I}\) 中的元素属于 \(\mathfrak{m}_{\mathbf{A}}\)，且有

\[
\dim_ {\mathrm{A}} (\mathrm{M}) = \operatorname{Card} (\mathrm{I}) + \dim_ {\mathrm{A}} \left(\mathrm{M} / \sum_ {i \in \mathrm{I}} x _ {i} \mathrm{M}\right),\tag{10}
\]

也就是说，若 \(i\mapsto x_{i}\) 是从 I 到 \(\mathfrak{m}_{\mathbf{A}}\) 中某个对于 M 为截切的子集的双射。

若 \(x\) 是 \(m_A\) 中的元素，称其为 \(M\) 的截切元，当且仅当 \(\{x\}\) 是 \(M\) 的截切子集，即有

\[
\dim_ {\mathbf {A}} (\mathbf {M} / x \mathbf {M}) = \dim_ {\mathbf {A}} (\mathbf {M}) - 1.
\]

备注。--- 1) 由公式 (6) 和 (7) 可知，S 是 M 的截切集，当且仅当它是 A/a 的截切集，其中 a 是 M 的湮灭子。

\begin{enumerate}
\def\labelenumi{\arabic{enumi})}
\setcounter{enumi}{1}
\tightlist
\item
  设 S 和 S′ 是 \(m_{A}\) 的两个不相交子集。要使 \(S \cup S'\) 成为 M 的截切集，必要且充分的条件是 S 为 M 的截切集，且 S' 为 \(M' = M/SM\) 的截切集。这由不等式 (8) 和公式
\end{enumerate}

\[
\begin{array}{r l} \operatorname{Card} (S \cup S ^ {\prime}) + \dim_ {A} (M / (S M + S ^ {\prime} M)) - \dim_ {A} (M) = & \\ = (\operatorname{Card} (S) + \dim_ {A} (M / S M) - \dim_ {A} (M)) + & \\ + (\operatorname{Card} (S ^ {\prime}) + \dim_ {A} (M ^ {\prime} / S ^ {\prime} M ^ {\prime}) - \dim_ {A} (M ^ {\prime})) \end{array}
\]

\begin{enumerate}
\def\labelenumi{\arabic{enumi})}
\setcounter{enumi}{2}
\tightlist
\item
  设 \(x \in \mathfrak{m}_{\mathbf{A}}\)，且 \(n \geqslant 1\) 为整数。显然，模 \(\mathbf{M} / x\mathbf{M}\) 和 \(\mathbf{M} / x^n\mathbf{M}\) 具有相同的支集，因而具有相同的维数。因此，\(x\) 是 \(\mathbf{M}\) 的截切元，当且仅当 \(x^n\) 是 \(\mathbf{M}\) 的截切元。由此以及备注 2，立即得到如下结果：设 \(x_1, ..., x_r\) 是 \(\mathfrak{m}_{\mathbf{A}}\) 中的元素，且 \(n_1, ..., n_r\) 是 \(> 0\) 的整数；则序列 \((x_1, ..., x_r)\) 是 \(\mathbf{M}\) 的截切序列，当且仅当序列 \((x_1^{n_1}, ..., x_r^{n_r})\) 是 \(\mathbf{M}\) 的截切序列。
\end{enumerate}

命题 3. --- 设 A 为诺特局部环，M 为非零有限生成 A-模。对于 \(\mathfrak{m}_{\mathrm{A}}\) 中的元素 x，x 是 M 的截切元的充要条件是：x 不属于任何极小元 \(\mathfrak{p}\)，其属于 \(\operatorname{Supp}(M)\) 且满足 \(\dim(\mathrm{A}/\mathfrak{p}) = \dim_{\mathrm{A}}(\mathrm{M})\)；而 M 上比率为 x 的倍乘映射 \(x_{\mathrm{M}}\) 为单射则是充分条件。

设 a 是 M 的湮灭子。说 x 是 M 的截切元，意味着 x 是 A/a 的截切元，而 M 的支集由 A 中满足 \(a \subset p\) 的素理想 p 构成；若 \(x_{M}\) 是单射，则 x 在 A/a 中的像不是 A/a 中的零因子。将 No. 1 的 Prop. 2 的 Cor. 2 应用于环 A/a，即得 Prop. 3。

推论。--- \(m_{A}\) 中对于 M 完全截切的每个元素序列（A, X, p. 157, def. 2）都是 M 的截切序列。

设 \((x_{1},\ldots ,x_{r})\) 是 \(\mathfrak{m}_{\mathbf{A}}\) 中对于 M 完全截切的元素序列。置 \(\mathbf{M}_0 = \mathbf{M}\)，并归纳定义 \(\mathbf{M}_i = \mathbf{M}_{i - 1} / x_i\mathbf{M}_{i - 1}\)，其中 \(1\leqslant i\leqslant r\)。根据 A, X, p. 160 的 cor. 1，比率为 \(x_{i}\) 且作用于 \(\mathbf{M}_{i - 1}\) 的倍乘映射是单射，因而 \(\dim_{\mathbf{A}}(\mathbf{M}_i) = \dim_{\mathbf{A}}(\mathbf{M}_{i - 1}) - 1\)（\(1\leqslant i\leqslant r\)）（prop. 3）。因此有

\[
\dim_ {\mathbf {A}} (\mathbf {M}) = r + \dim_ {\mathbf {A}} \left(\mathrm{M} / (x _ {1} \mathrm{M} + \dots + x _ {r} \mathrm{M})\right),
\]

所以序列 \((x_{1},...,x_{r})\) 是 M 的截切序列。

备注 4. --- 一般而言，M 的截切序列不一定是 M 的完全截切序列（p. 87, exerc. 6）。稍后我们将研究这样的诺特局部环上的模：其中每个截切序列都是完全截切的。

定理 1. --- 设 A 为诺特局部环，M 为非零有限生成 A-模，S 为 A 的极大理想 \(\mathfrak{m}_{\mathbf{A}}\) 的子集。

\begin{enumerate}
\def\labelenumi{\alph{enumi})}
\item
  若 M/SM 具有有限长度，则有 Card(S) \(\geqslant\) dim \(_{A}\) (M)；若 S 是 M 的截切集，则 Card(S) \(\leqslant\) dim \(_{A}\) (M)。
\item
  M 的每个截切子集都包含于 M 的某个极大截切子集中。
\item
  下列性质等价：
\end{enumerate}

\begin{enumerate}
\def\labelenumi{(\roman{enumi})}
\item
  S 是 M 的极大截切子集；
\item
  S 是 M 的截切子集，且 \(\operatorname{Card}(\mathbf{S}) = \dim_{\mathbf{A}}(\mathbf{M})\)；
\item
  M/SM 具有有限长度，且 Card(S) = dim\_A(M)；
\item
  S 是 M 的截切子集，且 M/SM 具有有限长度。
\end{enumerate}

由于 \(S \subset m_{A}\)，Nakayama 引理表明 M/SM ≠ \{0\}，因而 dim \(_{A}\) (M/SM) ≥ 0，且等号成立当且仅当 M/SM 具有有限长度。断言 \(a\) ) 随即由公式 (8)、(9) 以及性质 (ii)、(iii)、(iv) 的等价性得到。

断言 \(b\) ) 由如下事实得到：\(\mathfrak{m}_{\mathbf{A}}\) 中每个对于 \(\mathbf{M}\) 为截切的子集的基数都不超过整数 \(\dim_{\mathbf{A}}(\mathbf{M})\)。

由 \(a\) )，每个基数等于 \(\dim_{\mathbf{A}}(\mathbf{M})\) 的 M 的截切子集都是极大的。还需证明：若 S 是 M 的截切集且 \(\operatorname{Card}(\mathbf{S}) < \dim_{\mathbf{A}}(\mathbf{M})\)，则 S 不是极大的。设 a 是 M 的湮灭子，B 是局部诺特环 A/(a + SA)。根据 n° 1 的 prop. 2 的 cor. 2，存在元素 \(x\)，属于 \(\mathfrak{m}_{A}\)，使得 \(\dim(\mathrm{B}/x\mathrm{B}) = \dim(\mathrm{B}) - 1\)；因而 \(x \notin \mathrm{S}\)。根据备注 2，子集 S ∪ \(\{x\}\) 的 \(\mathfrak{m}_{A}\) 是 A/a 的截切集，由备注 1，它也是 M 的截切集。

推论。--- M 的维数是所有满足下述条件的整数 \(d \geqslant 0\) 中的最小者：存在一个序列 \((x_{1}, ..., x_{d})\)，其元素属于 \(m_{A}\)，使 A-模 M/ \(\sum_{i=1}^{d} x_{i}M\) 具有有限长度。

由于 \(\varnothing\) 是 M 的截切子集，定理 1 的 \(b\) ) 证明了 M 存在极大截切序列，即 \((x_{1},\ldots,x_{d})\)。有 \(d=\dim_{\mathbf{A}}(\mathbf{M})\) 且 A-模 \(\mathbf{M}/\sum_{i=1}^{d}x_{i}\mathbf{M}\) 具有有限长度，根据定理 1 的性质 (iii) 和 \(c\) )。反之，若 \((x_{1}',\ldots,x_{d'}')\) 是 \(\mathfrak{m}_{\mathbf{A}}\) 中元素组成的序列，使 A-模 \(\mathbf{M}/\sum_{j=1}^{d'}x_{j}'\mathbf{M}\) 具有有限长度，则有 \(d'\geqslant\dim_{\mathbf{A}}(\mathbf{M})\)，由 th. 1 的 \(a\) )。

回忆（III, § 3, n° 2, def. 1）：诺特局部环 A 的理想 q 称为 A 的定义理想，如果 A 的 q-进拓扑与 m\_A-进拓扑重合。

引理 2. --- 设 A 为诺特局部环，q 为 A 的理想。下列条件等价：

\begin{enumerate}
\def\labelenumi{(\roman{enumi})}
\item
  q 是 A 的定义理想；
\item
  存在整数 \(n \geqslant 0\)，使得 \(\mathfrak{m}_{\mathbf{A}}^{n} \subset \mathfrak{q} \subset \mathfrak{m}_{\mathbf{A}}\)；
\item
  q ≠ A，且 A/q 是具有有限长度的 A-模；
\item
  V(q) 等于 \(\{m_{A}\}\)（换言之，\(m_{A}\) 是包含 q 的唯一 A 的素理想）。
\end{enumerate}

事实上，(i)、(ii) 和 (iv) 的等价性已在 III, § 2, no. 5 中证明，而 (i) 与 (iii) 的等价性由 IV, § 2, no. 5, cor. 2 to prop. 9 得出。

因此，Thm. 1 的推论使我们可以陈述如下附论：

附论。--- 诺特局部环 A 的维数是满足存在一个由 d 个元素生成的 A 的定义理想的整数 \(d \geqslant 0\) 中的最小者。

